\documentclass[a4paper,12pt]{article}
\usepackage[utf8]{inputenc}
\usepackage[T1]{fontenc}
\usepackage[ngerman]{babel}
\usepackage{amsmath}
\usepackage{amssymb}
\usepackage{enumitem}
\usepackage{geometry}
\geometry{a4paper, margin=2.5cm}

\title{Einsendeaufgaben -- Lektion 2\\[0.5em]
\large Modul 61111: Mathematische Grundlagen}
\author{}
\date{}

\begin{document}
\maketitle

\section*{Aufgabe 2.6}

Sei $S \subseteq V$. Zu zeigen: $S$ ist Unterraum von $V$ $\iff$ $S = \langle S \rangle$.

\bigskip
\noindent\textbf{($\Rightarrow$)} Sei $S$ Unterraum von $V$. Zu zeigen: $S = \langle S \rangle$.

Per Kontraposition: Angenommen $S \neq \langle S \rangle$, dann existiert ein
$v \in \langle S \rangle$ mit $v \notin S$. Da $v$ als Linearkombination
\[
  v = \sum_{n} a_n x_n, \quad x_n \in S,
\]
dargestellt werden kann, aber $v \notin S$, ist $S$ nicht abgeschlossen unter
Linearkombinationen und daher kein Unterraum. \lightning

Die drei Unterraumeigenschaften zeigen $\langle S \rangle \subseteq S$ direkt:
\begin{enumerate}[label=\arabic*.]
  \item $0$ liegt immer in $\langle S \rangle$ (als leere Summe), also auch $0 \in S$.

  \item Sei $v, w \in S$. Wegen $S = \langle S \rangle$ folgt, dass $v + w$ als
  Linearkombination von Elementen aus $S$ dargestellt werden kann
  (mit Koeffizienten $1$ und $1$), also $v + w \in \langle S \rangle = S$.

  \item Sei $v \in S$ und $\lambda \in K$. Dann gilt:
  \[
    \lambda v = \lambda \sum_{n} a_n x_n = \sum_{n} (\lambda a_n) x_n
    \implies \lambda v \in \langle S \rangle \implies \lambda v \in S.
  \]
\end{enumerate}

\bigskip
\noindent\textbf{($\Leftarrow$)} Sei $S = \langle S \rangle$. Da $\langle S \rangle$
stets ein Unterraum ist, ist auch $S$ ein Unterraum.

\end{document}
